\chapter{RLC Circuits: Resonance, Damping, and Q}
\label{ch:rlc}
The RLC circuit is central to tuned receivers, oscillators, matching networks,
traps, crystal filters, and every selective RF stage. With two energy stores---a
capacitor and an inductor---it is the electrical embodiment of the
\emph{second-order system} of \cref{ch:linsys}, and everything that chapter said
about \(\omega_0\), \(\zeta\), damping regimes, pole geometry, and resonant peaks
applies here directly.

\section{The Circuit}
A source drives \(R\), \(L\), and \(C\) in series (\cref{fig:rlc-schematic}). Two
independent states carry the stored energy: the capacitor voltage \(v_C\) and the
inductor current \(i\).

\begin{figure}[htbp]
\centering
\begin{circuitikz}[scale=1]
\draw (0,0) to[sV,l=\(u(t)\)] (0,2)
      to[R,l=\(R\)] (2,2)
      to[L,l=\(L\),i>=\(i(t)\)] (4,2)
      to[C,l=\(C\),v=\(v_C\)] (4,0) -- (0,0);
\end{circuitikz}
\caption{Series RLC circuit. The two states are the capacitor voltage \(v_C\) and
the inductor current \(i\), so the system is second order.}
\label{fig:rlc-schematic}
\end{figure}

\section{From Schematic to Differential Equation}
KVL gives \(u=Ri+L\,\dd i/\dd t+v_C\) with \(i=C\,\dd v_C/\dd t\). Eliminating
\(i\) yields a second-order equation in the capacitor voltage:
\[
\boxed{\,LC\frac{\dd^2 v_C}{\dd t^2}+RC\frac{\dd v_C}{\dd t}+v_C=u(t).\,}
\]

\section{State-Space Form}
Using \(x_1=v_C\) and \(x_2=i\),
\[
\dot{\state}=
\mat{0 & 1/C\\ -1/L & -R/L}\state
+\mat{0\\ 1/L}u,
\qquad
y=v_C=\mat{1 & 0}\state .
\]
Two states, two poles---the eigenvalues of the state matrix.

\section{Matching the Generic Second-Order System}
Divide the ODE by \(LC\) to reach the standard form
\(\ddot v_C+2\zeta\omega_0\dot v_C+\omega_0^2 v_C=\omega_0^2 u\), giving the
identifications
\[
\omega_0=\frac{1}{\sqrt{LC}},
\qquad
\alpha=\frac{R}{2L},
\qquad
\zeta=\frac{\alpha}{\omega_0}=\frac{R}{2}\sqrt{\frac{C}{L}},
\qquad
Q=\frac{1}{2\zeta}=\frac{\omega_0 L}{R}.
\]
The poles are \(s=-\alpha\pm\jj\omega_d\) with \(\omega_d=\sqrt{\omega_0^2-\alpha^2}\).
Now the second-order results of \cref{ch:linsys} transfer wholesale: the step
responses of \cref{fig:step2} (over/critical/under-damped), the pole geometry of
\cref{fig:poles2} (on the \(\omega_0\) circle, \(\cos\theta=\zeta\)), and the
resonant peak of \cref{fig:bode2}. A tuned RF circuit is simply a very
lightly damped (high-\(Q\)) instance.

\section{Series Resonance}
The series impedance is
\[
Z=R+\jj\left(\omega L-\frac{1}{\omega C}\right).
\]
At \(\omega=\omega_0\) the reactances cancel, \(Z=R\) is purely resistive, and the
current is maximum for a fixed applied voltage. The reactive voltages across \(L\)
and \(C\) individually can then exceed the source voltage by a factor of \(Q\).

\section{Parallel Resonance}
The dual network---\(L\) and \(C\) in parallel---has admittance
\[
Y=\jj\omega C+\frac{1}{\jj\omega L}=\jj\left(\omega C-\frac{1}{\omega L}\right).
\]
At resonance the branch admittances cancel, so the input impedance is high
(infinite in the lossless ideal); real losses keep it finite. Series resonance is
a low-impedance notch in \(Z\); parallel resonance is a high-impedance peak.

\section{Q and Bandwidth}
The quality factor measures sharpness. For the series circuit
\(Q_s=\omega_0 L/R=1/(\omega_0 RC)\); for the parallel circuit the loss resistance
appears in parallel, so \(Q\) inverts to \(Q_p=R/(\omega_0 L)=\omega_0 RC\). In
both cases the half-power bandwidth is
\[
BW\approx\frac{f_0}{Q},
\]
and on the Bode magnitude (\cref{fig:bode2}) higher \(Q\) means a taller, narrower
peak---poles pressed against the \(\jj\omega\) axis.

\begin{mistakebox}
Do not treat one topology's \(Q\) formula as universal, and keep the three
\(Q\)'s straight. \emph{Loaded} \(Q\) includes source and load coupling;
\emph{unloaded} \(Q\) describes the resonator's own internal losses;
\emph{component} \(Q\) describes the loss of one inductor or capacitor. They
answer different questions.
\end{mistakebox}

\section{Root Locus: Damping Set by R}
Fix \(L\) and \(C\) and sweep \(R\). Since \(\zeta=\tfrac{R}{2}\sqrt{C/L}\), the
poles trace the canonical second-order locus of \cref{fig:rootlocus2}: at
\(R=0\) they sit on the \(\jj\omega\) axis at \(\pm\jj\omega_0\) (undamped,
infinite \(Q\)); as \(R\) grows they arc along the \(\omega_0\) circle into the
left half-plane, meet at \(-\omega_0\) when \(R=2\sqrt{L/C}\) (critical), then
split along the real axis (overdamped). Selectivity, ringing, and damping are the
one moving pole pair seen three ways.

\section{Energy Balance}
With \(\mathcal{H}=\tfrac12 Cv_C^2+\tfrac12 Li^2\), the series energy balance is
\[
\frac{\dd\mathcal{H}}{\dd t}=u(t)\,i(t)-R\,i^2(t).
\]
The inductor and capacitor trade energy back and forth at \(\omega_0\); the
resistor is the only drain. High \(Q\) means little is lost per cycle, so the
exchange rings for many cycles---the time-domain face of a sharp resonance.

\begin{controlsbox}
High \(Q\) \(\Leftrightarrow\) low \(\zeta\) \(\Leftrightarrow\) poles near the
\(\jj\omega\) axis \(\Leftrightarrow\) narrow bandwidth and long ringing. The tuned
circuit the exam calls ``selective'' is what a controls course calls ``lightly
damped.''
\end{controlsbox}

\section{Worked Example}
Let \(L=\SI{10}{\micro\henry}\), \(C=\SI{220}{\pico\farad}\), \(R=\SI{5}{\ohm}\).
Then
\[
f_0=\frac{1}{2\pi\sqrt{LC}}\approx\SI{3.39}{\mega\hertz},
\qquad
Q_s=\frac{\omega_0 L}{R}\approx 42.6,
\qquad
\zeta=\frac{1}{2Q}\approx0.0117.
\]
The poles sit at \(s=-\alpha\pm\jj\omega_d\approx-2.5\times10^5\pm\jj\,2.13\times10^7\ \si{\per\second}\):
the real part is about \(1/85\) of the imaginary part, so the poles are pressed
against the \(\jj\omega\) axis---a very high-\(Q\), sharply peaked circuit that
rings for roughly \(Q\approx43\) cycles. The half-power bandwidth is
\(BW\approx f_0/Q\approx\SI{80}{\kilo\hertz}\).
